\subsection{Objetivo}

Comparar la latencia de ida y vuelta (mediana y percentil~95) y los mensajes por segundo de sockets TCP con JSON y de gRPC con Protocol Buffers, en tres nodos que corren como procesos en \emph{localhost}, con mensajes de unos 100~bytes y de unos 10~KB y con uno o tres clientes enviando a la vez. Verificar, además, con los registros de los 20 intercambios, que los relojes de Lamport cumplen sus reglas en los dos mecanismos, y medir cómo cambia su tiempo de convergencia al pasar de 2 a 3 y a 4 nodos.

\subsection{Hipótesis}

Las hipótesis se escribieron antes de la corrida oficial. Suponen que cada cliente abre una sola conexión TCP o un solo canal de gRPC y hace por ahí sus 100 envíos, de modo que el saludo inicial de TCP y de HTTP/2 queda fuera de la medición. Los dos mecanismos corren en la misma máquina, con Python~3.13, y nunca a la vez.

En \emph{localhost} casi no hay red, por eso esperamos que la diferencia entre los dos mecanismos salga sobre todo del software que procesa cada mensaje. Con sockets, el cliente convierte el mensaje a JSON y lo envía con un salto de línea al final; el servidor lo lee, lo vuelve a convertir y responde con un acuse, también en JSON. Encima de TCP no hay otra capa. Con gRPC, cada llamada serializa el mensaje en binario, lo mete en tramas de HTTP/2 con sus encabezados y pasa por las capas de la biblioteca: los metadatos, el plazo de la llamada y el hilo del \texttt{ThreadPoolExecutor} que la atiende en el servidor.

Las dos fuentes que siguen no miden este mismo caso y se usan como evidencia indirecta. Zhu \emph{et al.}~\cite{Zhu2023} midieron el proxy Envoy, escrito en C++, con mensajes de 100~bytes: cuando procesaba el tráfico como gRPC añadía unas cinco veces más latencia que cuando lo trataba como TCP, y la mayor parte de ese costo extra venía de interpretar el protocolo. Es otro programa y otro lenguaje, pero muestra que las capas de gRPC tienen un costo por mensaje que TCP no paga. Niswar \emph{et al.}~\cite{Niswar2024} compararon gRPC con REST y con GraphQL, que usaban HTTP/1.1. gRPC tuvo menor tiempo de respuesta en la mayoría de sus pruebas, aunque en la de 100 solicitudes concurrentes con datos anidados REST fue más rápido. En ese estudio las otras opciones también cargan con HTTP. Aquí sockets no lo usa, por eso no esperamos que la ventaja de gRPC frente a REST se repita frente a sockets.

Con 100~bytes esperamos que pese sobre todo el costo fijo de cada llamada. Con 10~KB, el mensaje es en su mayor parte texto, y un texto ASCII ocupa casi lo mismo en los dos formatos: con nuestros campos, JSON suma unos 40~bytes más que Protocol Buffers por los nombres de los campos y las comillas. Eso es cerca del 40\,\% de un mensaje de 100~bytes y menos del 1\,\% de uno de 10~KB. Si el texto lleva tildes o eñes, JSON escribe cada una como \texttt{\textbackslash uXXXX} y la diferencia crece. Copiar y convertir esos 10~KB cuesta parecido en los dos mecanismos, de modo que la diferencia en milisegundos debería mantenerse y pesar menos frente al total.

Con tres clientes, los dos servidores atienden con hilos. En la versión estándar de CPython~3.13 solo un hilo a la vez ejecuta código Python, por el GIL. En gRPC, la red y HTTP/2 corren en el núcleo de la biblioteca, escrito en C, pero el método del servicio corre en un hilo de Python. No tenemos una fuente que diga cuál de los dos servidores sufre más esa espera. Suponemos que gRPC, porque cada llamada pasa también por código Python de la biblioteca antes y después del método. H3 pone a prueba esa suposición.

Los relojes de Lamport no llegan a un valor común, porque cada uno solo avanza. Por eso tomamos como convergencia el momento en que todos los nodos ya recibieron, directa o indirectamente, información de todos los demás, es decir, cuando en la historia causal de cada nodo hay al menos un evento de cada uno de los otros. El tiempo de convergencia va desde el primer envío de la corrida hasta la recepción con la que el último nodo cumple esa condición. Se calcula con la hora real de los registros, que se puede comparar entre nodos porque todos corren en la misma máquina y leen el mismo reloj. También contamos cuántos intercambios hicieron falta.

Los intercambios siguen una ronda en orden: el primer nodo envía al segundo, el segundo envía al tercero cuando recibe, y el último envía al primero. La ronda se repite hasta completar los 20 intercambios y sirve igual para 2, 3 y 4 nodos. Si se sigue mensaje por mensaje, todos los nodos saben de todos en el intercambio 2, 4 y 6, respectivamente, es decir, $2n-2$ para $n$ nodos. Como el patrón es fijo, ese número sale igual en todas las corridas, y las repeticiones solo aportan la variación del tiempo. Se hacen cinco corridas con cada cantidad de nodos y cada mecanismo.

\begin{description}
  \item[H1 (100~bytes, un cliente).] Sockets tiene menor mediana y menor percentil~95 de la latencia de ida y vuelta que gRPC, y más mensajes por segundo.
  \item[H2 (10~KB, un cliente).] Sockets sigue con menor mediana que gRPC, pero su ventaja relativa es menor que con 100~bytes. No esperamos que se invierta.
  \item[H3 (tres clientes a la vez, 100~bytes).] Frente a un cliente con 100~bytes, la ventaja relativa de sockets crece. Sockets mantiene menor mediana y menor percentil~95.
  \item[H4 (Lamport con 2, 3 y 4 nodos).] El tiempo de convergencia y el número de intercambios que necesita crecen al pasar de 2 a 3 y a 4 nodos, porque con más nodos la información de cada uno necesita más saltos para llegar a todos. Con la ronda esperamos 2, 4 y 6 intercambios.
\end{description}

\subsection{Criterio de decisión}

La unidad de análisis es cada repetición de 100 envíos, sin contar la de calentamiento. Por repetición, el script calcula la mediana y el percentil~95 de la latencia y los mensajes por segundo, que son los envíos de la repetición divididos entre su duración. Con tres clientes se juntan los 300 envíos de la repetición. Así quedan cinco valores por métrica, mecanismo y condición.

Hay diferencia real entre los dos mecanismos en una métrica solo si los rangos de sus cinco valores no se traslapan. Si se traslapan, no se puede afirmar que haya diferencia, y esa parte de la hipótesis no queda confirmada. H1 se cumple si las tres métricas muestran diferencia real a favor de sockets.

Para H2 y H3 no se emparejan repeticiones, así que el método sirve tanto si sockets y gRPC corren en bloques como si se intercalan. En cada condición se calculan los 25 cocientes posibles entre una de las cinco medianas de gRPC y una de las cinco de sockets. La ventaja relativa de sockets es menor con 10~KB si todos los cocientes con 10~KB quedan por debajo de todos los de 100~bytes, y crece con tres clientes si todos los cocientes con tres clientes quedan por encima de todos los de un cliente. Si los dos grupos de cocientes se traslapan, no se puede afirmar que la ventaja cambió.

H4 se cumple si, en cada mecanismo, los cinco tiempos de convergencia con 3 nodos quedan por encima de los cinco con 2, los de 4 por encima de los de 3, y el número de intercambios crece de 2 a 3 y a 4 nodos. Una corrida que no converge dentro de los 20 intercambios cuenta como mayor que cualquiera que sí convergió. Si el número de intercambios no da 2, 4 y 6, la ronda no se ejecutó como se describe aquí y hay que revisarla antes de discutir H4.
